\section{Conclusions}

In this project three distinct CNN architectures was proposed and trained: AlexNet, MobileNetV3, and SqueezeNet. Our analyses revealed that SqueezeNet generally outperformed the other models in key metrics like R2-Score, RMSE, and NRMSE. Concurrently, MobileNetV3 distinguished itself as the most efficient and fastest model, making it particularly suitable for deployment on resource-constrained hardware. Despite these promising outcomes, a significant limitation was identified: all models exhibited a notably high RMSE for ATE, approximately 1 meter. This level of error is likely too substantial for effectively correcting SLAM trajectories, and was primarily attributed to the presence of \quotes{broken} simulations where estimated SLAM trajectories were corrupted despite accurate final maps. 

For future works, the primary focus is to improve the dataset quality. It is crucial to obtain a larger dataset with significantly fewer \quotes{broken} simulations. Additionally, the proposed predictors could be employed to detect in real-time when SLAM errors exceed an acceptable threshold, allowing the system to be warned or, if necessary, to terminate the robot's run to prevent obtaining unusable maps.